\section{Model}
\label{sec:model}

\subsection{Timing, Assets, and Contracts}
\label{subsec:timing}

The model is static with two dates. At $t=0$, banks choose their funding mix; at $t=1$, all
contracts settle at face value. This two-date structure is the minimal environment for capturing
the funding-liquidity problem: institutions must commit to a portfolio before the full incidence
of dollar stress is known, and settlement reveals whether the choices were adequate.

\textbf{Assets.} There is a single safe collateral asset, US Treasury bonds, in fixed aggregate
supply $\bar{X} > 0$. Each Treasury bond has face value 1 (normalized) and trades at price
$P \in (0,1]$ at $t=0$. All dollar funding needs are denominated in US dollars.

\begin{quote}
\textbf{Assumption 0a} (First-order approximation). \textit{We normalize $P_0 = 1$ throughout.
This is valid to first order in $b^*$: since $P^* = 1 - \gamma\Theta b^*$ and typical basis
deviations satisfy $\gamma\Theta b^* \ll 1$ at calibrated parameter values, the error from
setting $P_0 = 1$ in bank-level decisions is of order $O(b^{*2})$. All formal results hold for
any $P_0 \in (0,1]$; the normalization $P_0 = 1$ simplifies notation.}
\end{quote}

Throughout the paper $P_0$ is retained symbolically in all formulas for generality; under
Assumption~0a, $\theta_i = P_0/\kappa_i = 1/\kappa_i$, $\Theta = 1/\kappa_2 + 1/\kappa_3$, and the
stability parameter simplifies as stated in \cref{subsec:equilibrium}.

\textbf{Contracts.} Two funding instruments are available to foreign banks at $t=0$.

\begin{enumerate}[label=(\roman*)]
  \item \textit{FX swap}: A contract in which the counterparty provides \$1 of dollar funding
    today against foreign-currency collateral at a pre-agreed exchange rate, with reversal at
    $t=1$. The cost per unit of swap is the \textit{cross-currency basis} $b \geq 0$, measured in
    annualized dollar terms. The basis is the price of dollar liquidity above the CIP-implied rate.

  \item \textit{Treasury sale}: A spot sale of Treasuries at market price $P$ per unit. A bank
    surrendering $x$ units receives $Px$ dollars. The liquidation incurs a convex cost
    $(\kappa_i/2)\,x^2$ reflecting bid-ask spreads, market impact, and roll-risk.
\end{enumerate}

\begin{quote}
\textbf{Assumption 0b} (Swap-only regime). \textit{In the main equilibrium, the dealer directs
all constrained balance-sheet capacity to swap supply, setting $Z^* = 0$. This is the conservative
bound: the dealer absorbs zero fire sales directly. Treasury fire sales are absorbed entirely by
outside investors. Appendix~A relaxes this assumption and verifies that all propositions are
robust.}
\end{quote}

\begin{quote}
\textbf{Assumption 0c} (Basel zero risk-weight). \textit{Treasury inventory carries zero
regulatory risk-weight, $q_T = 0$, consistent with the zero risk-weight assigned to US Treasuries
under Basel capital rules. Appendix~A shows that results extend to $q_T > 0$ by replacing $W_0$
with $W_0^{eff} \equiv W_0 - q_T\bar{X}_d/\phi$ throughout.}
\end{quote}

\subsection{The Backstop Hierarchy}
\label{subsec:hierarchy}

The central friction in the model is \textit{segmented access to official dollar liquidity}. This
is not a transaction cost or an informational asymmetry in the usual sense; it is a structural
feature of the international monetary architecture that differentiates US commercial banks, foreign
banks with swap-line access, and offshore dollar users along a single institutional dimension: the
coverage rate of their official backstop. The hierarchy maps directly to heterogeneous residual
funding needs, which in turn drive heterogeneous demand for market-based dollar liquidity — the
source of the equilibrium cross-currency basis.

This structure is the international analog of the inside/outside liquidity partition in
\citet{HolmstromTirole1998}. In their framework, inside liquidity (claims among private agents)
is insufficient to insure all aggregate shocks, motivating outside liquidity from the government.
Here, inside liquidity is private FX swap intermediation, and outside liquidity is the Federal
Reserve's dollar backstop. The three tiers formalize the degree to which each segment can access
outside liquidity.

\textbf{Segment 1 (US commercial banks, onshore).} These institutions have automatic,
unconditional access to Federal Reserve emergency facilities (discount window, FDIC guarantee).
Let $M_1 > 0$ denote segment 1's dollar funding need in the absence of backstop support. The
backstop fully covers this need:
\begin{equation}
  B_1 = M_1, \qquad R_1 \equiv M_1 - B_1 = 0.
  \label{eq:seg1}
\end{equation}
Segment 1 banks are inactive in equilibrium: they neither sell Treasuries nor enter FX swaps.
Their role in the model is to anchor the dollar at par and to fix the upper boundary of the
hierarchy. They do not solve an optimization problem.

\textbf{Segment 2 (foreign banks with swap lines).} These institutions belong to banking systems
whose central banks hold active bilateral dollar swap lines with the Federal Reserve. Let $M_2 > 0$
be the baseline dollar funding need of a representative segment 2 bank. The swap-line facility
covers fraction $\lambda \in [0,1]$ of this need:
\begin{equation}
  B_2 = \lambda M_2, \qquad R_2 \equiv M_2 - B_2 = (1-\lambda)M_2.
  \label{eq:seg2}
\end{equation}
The parameter $\lambda$ is the central policy variable of the model. At $\lambda = 1$, segment 2
is fully covered and its residual need vanishes; at $\lambda = 0$, it is equivalent to segment 3.
Segment 2 banks must cover $R_2$ using market instruments, subject to the optimization problem
stated in \cref{subsec:banks}.

\textbf{Segment 3 (offshore dollar users, no backstop).} These are foreign banks, investment
funds, and other financial entities with no access to any central-bank backstop facility. Let
$M_3 > 0$ be the baseline dollar funding need of a representative segment 3 entity. At $t=0$, a
funding shock $\varepsilon \geq 0$ hits this segment — an unexpected increase in dollar liabilities
or a sudden withdrawal of dollar credit lines. The backstop is zero:
\begin{equation}
  B_3 = 0, \qquad R_3 \equiv M_3 + \varepsilon.
  \label{eq:seg3}
\end{equation}
Segment 3 must cover $R_3$ entirely from market sources using the same instruments as segment 2,
with an optimization problem structurally identical to segment 2's (with $R_3$ replacing $R_2$).

The \textit{aggregate residual funding need} is the total dollar shortfall that must be met through
market instruments:
\begin{equation}
  R \equiv R_2 + R_3 = (1-\lambda)M_2 + M_3 + \varepsilon.
  \label{eq:R}
\end{equation}
The policy parameter $\lambda$ enters $R$ linearly with coefficient $-M_2 < 0$, establishing the
direct demand-relief channel of swap lines. The shock parameter $\varepsilon$ enters with
coefficient $+1$, representing exogenous dollar scarcity in the uncovered segment.

\subsection{Bank Optimization}
\label{subsec:banks}

Each bank in segment $i \in \{2,3\}$ with residual funding need $R_i > 0$ chooses the quantity
$x_i \geq 0$ of Treasuries to sell. The remaining need is covered by FX swaps:
\begin{equation}
  s_i = R_i - P x_i \geq 0. \tag{BC-$i$}
  \label{eq:BC}
\end{equation}
The bank faces a cost of FX swaps equal to $b$ per unit and a convex liquidation cost of $(\kappa_i/2)x_i^2$
on Treasury sales, capturing the private cost of unwinding positions rapidly. The optimization
problem is:
\begin{equation}
  \min_{x_i \geq 0} \quad b(R_i - P x_i) + \frac{\kappa_i}{2} x_i^2
  \qquad \text{subject to } x_i \leq R_i/P.
  \tag{P-$i$}
  \label{eq:Pi}
\end{equation}
Each bank takes the market price $P$ as given; $P$ is determined in equilibrium by market clearing.

The first-order condition for an interior solution (which holds when $b > 0$) is $-bP + \kappa_i x_i = 0$.
Evaluating at the reference price $P_0$, the optimal Treasury sales are:
\begin{equation}
  x_i^* = \theta_i b, \qquad \theta_i \equiv \frac{P_0}{\kappa_i}.
  \tag{FS-$i$}
  \label{eq:FSi}
\end{equation}
The scalar $\theta_i > 0$ is the \textit{basis-elasticity of fire sales} of segment $i$: a one-unit
increase in the basis induces additional Treasury sales of $\theta_i$ units. It is increasing in
$P_0$ — higher Treasury prices make liquidation more attractive — and decreasing in $\kappa_i$ —
higher liquidation costs deter sales. The quadratic cost is consistent with the fire-sale
friction in \citet{BrunnermeierPedersen2009}, adapted here to the international dollar-funding
context.

The aggregate basis-elasticity of Treasury fire sales across both segments is:
\begin{equation}
  \Theta \equiv \theta_2 + \theta_3 = P_0\!\left(\frac{1}{\kappa_2} + \frac{1}{\kappa_3}\right).
  \tag{$\Theta$-def}
  \label{eq:Theta}
\end{equation}
Substituting \eqref{eq:FSi} into \eqref{eq:BC} yields segment $i$'s swap demand:
\begin{equation}
  s_i^* = R_i - P_0\,\theta_i b = R_i - \frac{P_0^2}{\kappa_i}\,b.
  \tag{SD-$i$}
  \label{eq:SDi}
\end{equation}
Aggregating over segments 2 and 3:
\begin{equation}
  s_2^* + s_3^* = R - P_0\,\Theta b.
  \tag{SD-agg}
  \label{eq:SDagg}
\end{equation}
Aggregate swap demand is decreasing in the basis: as $b$ rises, banks substitute toward Treasury
sales, reducing swap demand. This downward-sloping demand is one of the two sides of the swap
market; the upward-sloping supply from the dealer is derived in \cref{subsec:dealer}.

\subsection{The Global Dealer}
\label{subsec:dealer}

There is a single representative global dealer — a large internationally active bank or primary
dealer that acts as market-maker in both the FX swap market and the Treasury secondary market.
The dealer is the sole supplier of FX swaps to segments 2 and 3. The dealer's balance sheet is
the mechanism through which fire sales feed back into swap supply.

\textbf{Balance sheet.} The dealer holds an initial Treasury inventory $\bar{X}_d > 0$ and initial
equity capital $W_0 > 0$. Its marked-to-market net worth at Treasury price $P$ is:
\begin{equation}
  W(P) = W_0 + P\bar{X}_d.
  \label{eq:wealth}
\end{equation}
The endogeneity of $W(P)$ is central to the model: when Treasury prices fall due to fire sales,
dealer wealth falls, tightening the leverage constraint and reducing swap supply. This is the
balance-sheet channel through which a segment 3 shock propagates to segment 2.

\textbf{Objective.} The dealer maximizes mean-variance utility over profits from FX swap
intermediation ($O$) and Treasury absorption ($Z$):
\begin{align}
  \max_{O \geq 0,\, Z \geq 0} \quad bO + (1-P)Z - \frac{\eta}{2}\!\left(\sigma_O^2 O^2 + \sigma_Z^2 Z^2\right),
  \tag{P-D}
  \label{eq:PD}
\end{align}
where $\eta > 0$ is the CARA risk-aversion coefficient, $\sigma_O^2$ is the variance of payoff per
unit of swap position, and $\sigma_Z^2$ is the variance per unit of Treasury absorption. The objective
is strictly concave.

\textbf{Leverage constraint.} The dealer is subject to a Basel-type regulatory leverage constraint:
\begin{equation}
  q_S O + q_T(\bar{X}_d + Z) \leq \phi W(P),
  \tag{LC}
  \label{eq:LC}
\end{equation}
where $q_S > 0$ and $q_T \geq 0$ are the risk-weights on swap and Treasury positions respectively,
and $\phi > 0$ is the maximum permitted leverage ratio. The feasible set is convex and compact for
fixed $P$; a unique solution to \eqref{eq:PD} subject to \eqref{eq:LC} exists.

\textit{Remark} (Parameter roles). The risk-aversion parameters $(\eta, \sigma_O^2, \sigma_Z^2)$
govern the unconstrained dealer interior solution. In the constrained regime — the empirically
relevant case emphasized by \citet{Duffie2020}, \citet{DuTepperVerdelhan2018}, and
\citet{GabaixMaggori2015} — the leverage constraint~\eqref{eq:LC} binds and swap supply is
$O^* = \Gamma W(P)$ regardless of $(\eta, \sigma_O^2, \sigma_Z^2)$. These parameters do not
appear in the main propositions; their role is confined to determining the threshold wealth level
below which \eqref{eq:LC} binds, which is characterized in Appendix~A.

\textbf{Constrained equilibrium.} The model focuses on the regime in which~\eqref{eq:LC} binds.
Under Assumptions~0b and~0c ($Z^* = 0$, $q_T = 0$), the binding constraint reduces to
$q_S O \leq \phi W(P)$, yielding:
\begin{equation}
  O^* = \Gamma W(P), \qquad \Gamma \equiv \frac{\phi}{q_S}.
  \tag{O-star}
  \label{eq:Ostar}
\end{equation}
The parameter $\Gamma = \phi/q_S > 0$ is the \textit{dealer capacity parameter}: it is increasing
in the regulatory leverage limit $\phi$ and decreasing in the swap risk-weight $q_S$. This
directly links the architecture of \citet{GabaixMaggori2015} — constrained financier capacity
governing the price of intermediation — to the specific Basel parameters of a dollar swap dealer.
Under the endogenous wealth equation~\eqref{eq:wealth}, swap supply depends on the Treasury price:
when fire sales depress $P$, dealer wealth contracts, leverage tightens, and swap supply falls.
This is the amplification channel that Propositions 2–4 characterize in closed form.

\subsection{Outside Investors and the Treasury Market}
\label{subsec:outsideinvestors}

A continuum of outside investors — pension funds, sovereign wealth funds, reserve managers —
supplies a downward-sloping demand curve for Treasuries. Each investor has CARA preferences with
absolute risk aversion $\alpha > 0$ and forms the posterior that each Treasury pays face value 1
at $t=1$ with variance $\sigma_T^2$. The standard mean-variance portfolio problem yields aggregate
outside-investor demand:
\begin{equation}
  X^{out}(P) = \frac{1-P}{\gamma}, \qquad \gamma \equiv \alpha\sigma_T^2 > 0.
  \tag{OI-dem}
  \label{eq:OIdem}
\end{equation}
The baseline demand at $P = 1$ is normalized to zero; the elasticity parameter $\gamma$ is
the price sensitivity of outside investors to departures from par. Under Assumption~0b, outside
investors absorb all bank fire sales: the dealer does not purchase Treasuries in equilibrium, so
$X^{out}(P^*)$ must equal total bank sales $x_2^* + x_3^*$.

Inverting~\eqref{eq:OIdem} at the clearing quantity gives the equilibrium price equation:
\begin{equation}
  P^* = 1 - \gamma\Theta b.
  \tag{P-eq}
  \label{eq:Peq}
\end{equation}
A higher basis induces more fire sales, pushing down the Treasury price. The price-impact
elasticity is $\gamma\Theta$, the product of the outside-investor sensitivity $\gamma$ and the
aggregate fire-sale elasticity $\Theta$. These two objects are never conflated in the derivations
that follow.

\subsection{Equilibrium}
\label{subsec:equilibrium}

\begin{definition}[Static Equilibrium]
\label{def:equilibrium}
A \textit{static equilibrium} is a tuple $(b^*, P^*, O^*, Z^*, \{x_i^*, s_i^*\}_{i=2,3})$
satisfying:
\begin{enumerate}[label=(\roman*)]
  \item \textbf{Bank optimization:} For each $i \in \{2,3\}$, $(x_i^*, s_i^*)$ solves
    \emph{(P-$i$)} given $b^*$ and $P^*$.
  \item \textbf{Dealer optimization:} $(O^*, Z^*)$ solves \emph{(P-D)} subject to \emph{(LC)}
    given $b^*$ and $P^*$, with the leverage constraint binding.
  \item \textbf{FX swap market clearing:}
    \begin{equation}
      s_2^* + s_3^* = O^*. \tag{MC-S} \label{eq:MCS}
    \end{equation}
  \item \textbf{Treasury market clearing} (under Assumption~0b):
    \begin{equation}
      x_2^* + x_3^* = X^{out}(P^*) = \frac{1-P^*}{\gamma}. \tag{MC-T} \label{eq:MCT}
    \end{equation}
  \item \textbf{Non-negativity:} $b^* \geq 0$, $P^* \in (0,1]$, $x_i^* \geq 0$, $s_i^* \geq 0$,
    $O^* \geq 0$, $Z^* = 0$.
\end{enumerate}
\end{definition}

\textbf{Reduction to a single equation.} The equilibrium conditions collapse to a single scalar
equation in $b^*$ through four steps.

\textit{Step 1: Treasury price.} From~\eqref{eq:FSi}, aggregate Treasury sales are $x_2^* + x_3^* = \Theta b$.
Equating with outside-investor demand via~\eqref{eq:MCT} and~\eqref{eq:OIdem} and solving for $P$
gives \eqref{eq:Peq}: $P^* = 1 - \gamma\Theta b$.

\textit{Step 2: Dealer wealth.} Substituting~\eqref{eq:Peq} into~\eqref{eq:wealth}:
\begin{equation}
  W(b) = W_0 + (1 - \gamma\Theta b)\bar{X}_d = \tilde{W} - \gamma\Theta b\,\bar{X}_d,
  \qquad \tilde{W} \equiv W_0 + \bar{X}_d.
  \tag{W-b}
  \label{eq:Wb}
\end{equation}
Dealer wealth falls linearly with the basis: $\partial W/\partial b = -\gamma\Theta\bar{X}_d < 0$.

\textit{Step 3: Swap supply.} From~\eqref{eq:Ostar} and~\eqref{eq:Wb}:
\begin{equation}
  O^*(b) = \Gamma\!\left(\tilde{W} - \gamma\Theta b\,\bar{X}_d\right).
  \tag{O-b}
  \label{eq:Ob}
\end{equation}
Swap supply is decreasing in $b$: a higher basis is associated with lower Treasury prices, lower
dealer wealth, and lower swap capacity. This decreasing supply, combined with the decreasing
aggregate demand~\eqref{eq:SDagg}, is what makes the equilibrium condition well-defined.

\textit{Step 4: FX swap market clearing.} Equating~\eqref{eq:SDagg} with~\eqref{eq:Ob}:
\[
  R - P_0\Theta b = \Gamma\!\left(\tilde{W} - \gamma\Theta b\,\bar{X}_d\right).
\]
Collecting terms in $b$ and factoring:
\[
  R - \Gamma\tilde{W} = b\Theta(P_0 - \Gamma\gamma\bar{X}_d).
\]

\begin{definition}[Stability Parameter]
\label{def:delta}
\begin{equation}
  \Delta \equiv \Theta(P_0 - \Gamma\gamma\bar{X}_d).
  \tag{$\Delta$-def}
  \label{eq:Delta}
\end{equation}
\end{definition}

The equilibrium basis is:
\begin{equation}
  b^* = \frac{R - \Gamma\tilde{W}}{\Delta}.
  \tag{EQ}
  \label{eq:EQ}
\end{equation}

Two conditions ensure a positive, well-defined equilibrium.

\begin{quote}
\textbf{Assumption 1} (Stability). $\Delta > 0$, equivalently $P_0 > \Gamma\gamma\bar{X}_d$.
\textit{This requires the direct demand-relief effect of a higher basis to dominate the
amplification effect through dealer wealth. When $\Delta \leq 0$, the feedback loop overwhelms
the stabilizing mechanism and no finite equilibrium basis clears the market.}
\end{quote}

\begin{quote}
\textbf{Assumption 2} (Positive Basis). $R > \Gamma\tilde{W}$.
\textit{The aggregate funding need exceeds the dealer's unconstrained swap capacity at $b = 0$.
This is the condition for dollar scarcity: without it, the dealer absorbs all required swaps at
zero basis premium.}
\end{quote}

\begin{proposition}[Existence and Uniqueness]
\label{prop:existence}
Under Assumptions~1 and~2, there exists a unique equilibrium basis $b^* > 0$ given by
equation~\eqref{eq:EQ}. The corresponding equilibrium Treasury price is
$P^* = 1 - \gamma\Theta b^* \in (0,1)$, and all equilibrium quantities are strictly positive.
\end{proposition}

\begin{proof}
Under Assumption~1, $\Delta = \Theta(P_0 - \Gamma\gamma\bar{X}_d) > 0$, so the denominator of
\eqref{eq:EQ} is strictly positive. Under Assumption~2, the numerator $R - \Gamma\tilde{W} > 0$.
Therefore $b^* = (R - \Gamma\tilde{W})/\Delta > 0$ is well-defined and unique.

For the Treasury price: $P^* = 1 - \gamma\Theta b^* < 1$ since $b^* > 0$. That $P^* > 0$ requires
$\gamma\Theta b^* < 1$, i.e., $(R - \Gamma\tilde{W})/(\Delta\gamma\Theta) < 1/(\gamma\Theta)$,
which reduces to $R - \Gamma\tilde{W} < \Delta = \Theta(P_0 - \Gamma\gamma\bar{X}_d)$, i.e.,
$R < P_0\Theta + \Gamma W_0$. This upper bound on the size of the funding shock is maintained
as a parametric restriction without requiring a separate assumption.

Swap demands $s_i^* = R_i - P_0\theta_i b^* > 0$ for each $i \in \{2,3\}$ under the same upper
bound. Fire sales $x_i^* = \theta_i b^* > 0$ under Assumption~2. Dealer supply
$O^* = \Gamma W(b^*) > 0$ since $W(b^*) = \tilde{W} - \gamma\Theta b^*\bar{X}_d > 0$
by $P^* > 0$.

All equilibrium conditions (i)--(v) of \cref{def:equilibrium} are satisfied by construction.
\end{proof}

\Cref{prop:existence} establishes that the backstop hierarchy, the dealer's
leverage constraint, and the fire-sale dynamics together produce a unique price of dollar
liquidity. The equilibrium basis $b^*$ in~\eqref{eq:EQ} is a closed-form function of the
structural primitives: the aggregate funding need $R$ (which depends on the backstop parameter
$\lambda$ and the stress shock $\varepsilon$), the dealer capacity $\Gamma$, the dealer wealth at par
$\tilde{W}$, the aggregate fire-sale elasticity $\Theta$, the outside-investor elasticity $\gamma$,
and the dealer Treasury inventory $\bar{X}_d$. The amplification properties of this equilibrium
are the subject of \cref{sec:results}.
